\documentclass[../../main.tex]{subfiles}
\begin{document}

{\bf [解]}
\textbf{(1)} 
$t$によって切断面の方程式が変化するので，これを
\begin{align}
 \pi(t): x+\dfrac{1}{2}y+\dfrac{1}{3}z=t \label{eq:1}
\end{align}
とおく．

下図のように直方体$V$の頂点を定める．O,A,B,Cは$y=0$上にあり，D,E,F,Gは$y=2$上に存在する．
この時，直方体$V$の体積は
\begin{align}
 T = 6 \label{eq:2}
\end{align}
である．以下$t$の大きさで場合分けする．

% 視点の設定 (tdplot_set_coordinates{r}{theta}{phi})
% 手書きスケッチに近い視点 (右上から見下ろす)
\tdplotsetmaincoords{65}{120}

\begin{figure}[htb]
\begin{center}
\begin{tikzpicture}[tdplot_main_coords, scale=1.5]

    % --- 1. 座標定義 ---
    % 直方体 V の頂点 (0<=x<=1, 0<=y<=2, 0<=z<=3)
    % 手書き図の対応:
    % y=0 上: O(0,0,0), A(0,0,3), B(1,0,3), C(1,0,0)
    % y=2 上: G(0,2,0), D(0,2,3), E(1,2,3), F(1,2,0)
    \coordinate (O) at (0,0,0);
    \coordinate (A) at (0,0,3);
    \coordinate (B) at (1,0,3);
    \coordinate (C) at (1,0,0);
    \coordinate (G) at (0,2,0);
    \coordinate (D) at (0,2,3);
    \coordinate (E) at (1,2,3);
    \coordinate (F) at (1,2,0);

    % --- 2. 座標軸の描画 ---
    \draw[->, thick, gray] (0,0,0) -- (2.0,0,0) node[anchor=north east]{$x$};
    \draw[->, thick, gray] (0,0,0) -- (0,3.0,0) node[anchor=north west]{$y$};
    \draw[->, thick, gray] (0,0,0) -- (0,0,4.0) node[anchor=south]{$z$};

    % --- 3. 切平面 pi(t) と軸の交点 (t = 0.8 の例) ---
    % x + y/2 + z/3 = t  =>  x-int=t, y-int=2t, z-int=3t
    \pgfmathsetmacro{\tval}{0.8}
    \coordinate (Px) at (\tval, 0, 0);
    \coordinate (Py) at (0, 2*\tval, 0);
    \coordinate (Pz) at (0, 0, 3*\tval);

    % 切断領域 (三角錐) の塗りつぶし
    \fill[red!30, opacity=0.6] (Px) -- (Py) -- (Pz) -- cycle;
    \draw[red!80!black, thick] (Px) -- (Py) -- (Pz) -- cycle;

    % --- 4. 直方体 V の辺の描画 ---
    % 隠線 (破線)
    \draw[dashed, thick] (O) -- (C);
    \draw[dashed, thick] (O) -- (G);
    \draw[dashed, thick] (O) -- (A);

    % 実線
    \draw[thick] (C) -- (B) -- (A);
    \draw[thick] (G) -- (F) -- (E) -- (D) -- cycle;
    \draw[thick] (C) -- (F);
    \draw[thick] (B) -- (E);
    \draw[thick] (A) -- (D);

    % --- 5. 目盛りラベル ---
    \node[anchor=north east] at (C) {$1$};
    \node[anchor=south west] at (G) {$2$};
    \node[anchor=west] at (A) {$3$};

    % --- 6. 頂点ラベルと点のプロット ---
    \fill[black] (O) circle (1pt) node[anchor=north west] {$O$};
    \fill[black] (A) circle (1pt) node[anchor=south east] {$A$};
    \fill[black] (B) circle (1pt) node[anchor=north west] {$B$};
    \fill[black] (C) circle (1pt) node[anchor=south east] {$C$};
    \fill[black] (D) circle (1pt) node[anchor=south east] {$D$};
    \fill[black] (E) circle (1pt) node[anchor=south west] {$E$};
    \fill[black] (F) circle (1pt) node[anchor=west] {$F$};
    \fill[black] (G) circle (1pt) node[anchor=south east] {$G$};

    \node[above right] at (E) {$V$};

\end{tikzpicture}
\caption{直方体Vの様子．}
\end{center} 
\end{figure}


\subparagraph{$1^\circ$ $0<t\le1$の時}

この時切断面は以下のようになり，直方体の辺と切断面の交点は辺OC, OG, OA上にくる．

% 視点の設定 (tdplot_set_coordinates{r}{theta}{phi})
\tdplotsetmaincoords{65}{120}

\begin{figure}[htb]
\begin{center}
\begin{tikzpicture}[tdplot_main_coords, scale=1.5]

    % --- 1. 座標定義 ---
    % 直方体 V の頂点
    \coordinate (O) at (0,0,0);
    \coordinate (A) at (0,0,3);
    \coordinate (B) at (1,0,3);
    \coordinate (C) at (1,0,0);
    \coordinate (G) at (0,2,0);
    \coordinate (D) at (0,2,3);
    \coordinate (E) at (1,2,3);
    \coordinate (F) at (1,2,0);

    % --- 2. 座標軸の描画 ---
    \draw[->, thick, gray] (0,0,0) -- (2.2,0,0) node[anchor=north east]{$x$};
    \draw[->, thick, gray] (0,0,0) -- (0,3.2,0) node[anchor=north west]{$y$};
    \draw[->, thick, gray] (0,0,0) -- (0,0,4.2) node[anchor=south]{$z$};

    % --- 3. 切平面 pi(t) と軸の交点 (0 < t <= 1) ---
    % x-int = t, y-int = 2t, z-int = 3t
    \pgfmathsetmacro{\tval}{0.7}
    \coordinate (Px) at (\tval, 0, 0);       % x軸上の点 (t)
    \coordinate (Py) at (0, 2*\tval, 0);     % y軸上の点 (2t)
    \coordinate (Pz) at (0, 0, 3*\tval);     % z軸上の点 (3t)

    % 切り口 (三角形) の塗りつぶしと境界線
    \fill[red!25, opacity=0.6] (Px) -- (Py) -- (Pz) -- cycle;
    \draw[red!80!black, ultra thick] (Px) -- (Py) -- (Pz) -- cycle;

    % --- 4. 直方体 V のフレーム描画 ---
    % 隠線 (破線)
    \draw[dashed, thick] (O) -- (C);
    \draw[dashed, thick] (O) -- (G);
    \draw[dashed, thick] (O) -- (A);

    % 外枠 (実線)
    \draw[thick] (C) -- (B) -- (A);
    \draw[thick] (G) -- (F) -- (E) -- (D) -- cycle;
    \draw[thick] (C) -- (F);
    \draw[thick] (B) -- (E);
    \draw[thick] (A) -- (D);

    % --- 5. 交点 (t, 2t, 3t) のプロットとラベル ---
    \fill[red!80!black] (Px) circle (1.2pt) node[anchor=north west] {$t$};
    \fill[red!80!black] (Py) circle (1.2pt) node[anchor=south west] {$2t$};
    \fill[red!80!black] (Pz) circle (1.2pt) node[anchor=east] {$3t$};

    % --- 6. 頂点ラベル ---
    \fill[black] (O) circle (1pt) node[anchor=north east] {$O$};
    \fill[black] (A) circle (1pt) node[anchor=south east] {$A$};
    \fill[black] (B) circle (1pt) node[anchor=north west] {$B$};
    \fill[black] (C) circle (1pt) node[anchor=south west] {$C$};
    \fill[black] (D) circle (1pt) node[anchor=south east] {$D$};
    \fill[black] (E) circle (1pt) node[anchor=south west] {$E$};
    \fill[black] (F) circle (1pt) node[anchor=west] {$F$};
    \fill[black] (G) circle (1pt) node[anchor=south east] {$G$};
\end{tikzpicture}
\caption{$0<t\le 1$の時の切断面}
\end{center}
\end{figure}

切断された立体のうち点Oを含む三角錐の体積は
\begin{align}
\dfrac{1}{3}\cdot t\cdot\left(\dfrac{1}{2}\cdot t\cdot 2t\cdot\dfrac{3}{2}\right)=t^3
\end{align}
であり，これは$\dfrac{T}{2}=3$よりも小さいから題意より
\begin{align}
 f(t) = t^3
\end{align}
である．

\subparagraph{$2^\circ$ $1\le t\le2$の時}

この時切断面は以下のようになり，直方体の辺と切断面の交点は辺BC,CF, FG, GD, DA, AB上にくる．

% 視点の設定 (tdplot_set_coordinates{r}{theta}{phi})
\tdplotsetmaincoords{65}{120}

\begin{figure}[htb]
 \begin{center}
\begin{tikzpicture}[tdplot_main_coords, scale=1.5]

    % --- 1. 座標定義 ---
    % 直方体 V の頂点
    \coordinate (O) at (0,0,0);
    \coordinate (A) at (0,0,3);
    \coordinate (B) at (1,0,3);
    \coordinate (C) at (1,0,0);
    \coordinate (G) at (0,2,0);
    \coordinate (D) at (0,2,3);
    \coordinate (E) at (1,2,3);
    \coordinate (F) at (1,2,0);

    % パラメータ設定 (1 < t < 1.5 例として t = 1.25)
    \pgfmathsetmacro{\tval}{1.25}

    % 大きな三角錐の頂点 (各軸との交点 t, 2t, 3t)
    \coordinate (Px) at (\tval, 0, 0);       % x軸上: t
    \coordinate (Py) at (0, 2*\tval, 0);     % y軸上: 2t
    \coordinate (Pz) at (0, 0, 3*\tval);     % z軸上: 3t

    % 切断六角形の頂点 (直方体の各辺との交点)
    \coordinate (K_CF) at (1, {2*(\tval-1)}, 0);    % 辺 CF 上
    \coordinate (K_FG) at ({\tval-1}, 2, 0);        % 辺 FG 上
    \coordinate (K_GD) at (0, 2, {3*(\tval-1)});    % 辺 GD 上
    \coordinate (K_DA) at (0, {2*(\tval-1)}, 3);    % 辺 DA 上
    \coordinate (K_AB) at ({\tval-1}, 0, 3);        % 辺 AB 上
    \coordinate (K_BC) at (1, 0, {3*(\tval-1)});    % 辺 BC 上

    % --- 2. 座標軸の描画 ---
    \draw[->, thick, gray] (0,0,0) -- (2.0,0,0) node[anchor=north east]{$x$};
    \draw[->, thick, gray] (0,0,0) -- (0,3.2,0) node[anchor=north west]{$y$};
    \draw[->, thick, gray] (0,0,0) -- (0,0,4.2) node[anchor=south]{$z$};

    % --- 3. 大きな三角錐の破線外枠 (手書き図の t, 2t, 3t を結ぶ直線) ---
    \draw[dashed, red!60!black, thin] (Px) -- (Py) -- (Pz) -- cycle;

    % --- 4. 切り口 (六角形) の塗りつぶしと境界線 ---
    \fill[blue!25, opacity=0.6] 
        (K_CF) -- (K_FG) -- (K_GD) -- (K_DA) -- (K_AB) -- (K_BC) -- cycle;
    
    \draw[blue!80!black, ultra thick] 
        (K_CF) -- (K_FG) -- (K_GD) -- (K_DA) -- (K_AB) -- (K_BC) -- cycle;

    % --- 5. 直方体 V のフレーム描画 ---
    % 隠線 (破線)
    \draw[dashed, thick] (O) -- (C);
    \draw[dashed, thick] (O) -- (G);
    \draw[dashed, thick] (O) -- (A);

    % 外枠 (実線)
    \draw[thick] (C) -- (B) -- (A);
    \draw[thick] (G) -- (F) -- (E) -- (D) -- cycle;
    \draw[thick] (C) -- (F);
    \draw[thick] (B) -- (E);
    \draw[thick] (A) -- (D);

    % --- 6. 各辺上の切断点 (黒丸) ---
    \fill[blue!80!black] (K_CF) circle (1.3pt);
    \fill[blue!80!black] (K_FG) circle (1.3pt);
    \fill[blue!80!black] (K_GD) circle (1.3pt);
    \fill[blue!80!black] (K_DA) circle (1.3pt);
    \fill[blue!80!black] (K_AB) circle (1.3pt);
    \fill[blue!80!black] (K_BC) circle (1.3pt);

    % --- 7. 交点 (t, 2t, 3t) のプロットとラベル ---
    \fill[red!60!black] (Px) circle (1pt) node[anchor=north west] {$t$};
    \fill[red!60!black] (Py) circle (1pt) node[anchor=south west] {$2t$};
    \fill[red!60!black] (Pz) circle (1pt) node[anchor=east] {$3t$};

    % --- 8. 頂点ラベル ---
    \fill[black] (O) circle (1pt) node[anchor=north east] {$O$};
    \fill[black] (A) circle (1pt) node[anchor=south east] {$A$};
    \fill[black] (B) circle (1pt) node[anchor=north west] {$B$};
    \fill[black] (C) circle (1pt) node[anchor=south west] {$C$};
    \fill[black] (D) circle (1pt) node[anchor=south east] {$D$};
    \fill[black] (E) circle (1pt) node[anchor=south west] {$E$};
    \fill[black] (F) circle (1pt) node[anchor=west] {$F$};
    \fill[black] (G) circle (1pt) node[anchor=south east] {$G$};
\end{tikzpicture}
\caption{$1\le t \le 2$の時の切断面．}
\end{center}
\end{figure}

切断立体のうち原点$O$を含む方の体積は，大きな三角錐からコーナーの三角錐を3つ分切り取ることで求められる．
大きな三角錐の体積は
\begin{align}
 \dfrac{1}{3}2t \dfrac{t\cdot 3t}{2} = t^3
\end{align}
である．対してコーナーの小さい三角錐は大きな三角錐との比が$t:t-1$, $2t:2t-2$,$3t:3t-3$でいずれも$t:t-1$となるから，いずれも体積は
\begin{align}
 t^3\cdot \dfrac{(t-1)^3}{t^3} = (t-1)^3
\end{align}
である．従って断立体のうち原点$O$を含む方の体積は
\begin{align}
  t^3-3(t-1)^3=-2t^3+9t^2-9t+3 \equiv g(t)
\end{align}
である．$g'(t)=-6(t^2-3t+2)=-6(t-2)(t-1)\ge0\ (1\le t\le2)$だから$g(t)$は非減少で，かつ$g(3/2)=3$だから，
\begin{align}
f(t)=
\begin{cases}\displaystyle
g(t) & \left(1\le t\le\dfrac{3}{2}\right) \\
6-g(t) & \left(\dfrac{3}{2}\le t\le2\right)
\end{cases}
\end{align}

\subparagraph{$3^\circ$ $2\le t<3$の時}

$t$についての対称性から $f(t)=(3-t)^3$ である．

以上3つの場合分けから求める体積は
\begin{align}
f(t)=
\begin{cases}\displaystyle
t^3 & (0\le t\le1) \\
-2t^3+9t^2-9t+3 & \left(1\le t\le\dfrac{3}{2}\right) \\
2t^3-9t^2+9t+3 & \left(\dfrac{3}{2}\le t\le2\right) \\
(3-t)^3 & (2\le t\le3)
\end{cases}
\end{align}
である．


(2)

$f(t)$は全区間で連続だから題意を示すには$t=1,\dfrac{3}{2}$での両側極限の一致を見れば良い．
簡単のため
\begin{align}
 p(t) = t^3 \\
 q(t) = 6-g(t)
\end{align}
と置くと，これらは$t=1,\dfrac{3}{2}$で微分可能である．

\subparagraph{1. $t=1$の時}

この時の左右極限は
\begin{align}
 \lim_{h\to +0}\dfrac{f(1+h)-f(1)}{h} = p'(1)=3, \qquad \lim_{h\to +0}\dfrac{f(1-h)-f(1)}{-h} = g'(1)=3
\end{align}
より一致するから，$f(t)$は$t=1$で微分可能である．

\subparagraph{2. $t=\dfrac{3}{2}$の時}

この時の左右極限は
\begin{align}
  \lim_{h\to +0}\dfrac{f(3/2+h)-f(3/2)}{h} = q'(3/2)=-\dfrac{9}{2}, \qquad \lim_{h\to +0}\dfrac{f(3/2-h)-f(3/2)}{-h} = g'(3/2)=+\dfrac{9}{2}
\end{align}
より一致しないから，$f(t)$は$t=\dfrac{3}{2}$では微分不可能である．

以上から
\begin{align}
t=1:\ \text{微分可能}, \qquad t=\dfrac{3}{2}:\ \text{微分不可能}
\end{align}
である．


\newpage

\end{document}
